\chapter{Smooth fields and formal Hamiltonians}
\label{chap:smooth-formal-hamiltonians}

\section{Holomorphic fields and Hamiltonian parameters}
\label{sec:tb-fields}

A local expression involving finitely many derivatives of a smooth field can
be evaluated at a point without requiring an analytic expansion there.
Formal Hamiltonian algebra uses exactly such derivative data, assembled into
a power series.  The distinction becomes visible in one variable:
\[
 F(x,y)=\sum_{n\ge1}n!x^n.
\]
Every coefficient can occur as a holomorphic derivative of a smooth function.
No holomorphic function near the origin has this Taylor series.  The
comparison between smooth and formal fields must therefore preserve the
differential that imposes holomorphicity.

Throughout, \(\mathbb N=\{0,1,2,\ldots\}\), and all vector spaces and
multilinear maps are over \(\mathbb C\), unless a real coordinate is
specified.

The Dolbeault differential supplies this differential.  In a holomorphic
field description, degree-zero functions give Hamiltonian parameters and
their antiholomorphic derivatives measure failure of holomorphicity.
The Poisson bracket combines their infinitesimal transformations.
Scalar cochains then encode alternating functions of these symmetry
parameters.  The construction below concerns this symmetry algebra and its
scalar cochains.  It assumes no action functional, integration cycle, or
equivalence of physical theories.

Fix an open polydisc
\[
 U=D(0,R_1)\times D(0,R_2)\subset\mathbb C^2,
 \qquad 0<R_1,R_2<\infty.
\]
Here \(D(a,r)=\{z\in\mathbb C:|z-a|<r\}\).
Write \(z_j=s_j+it_j\), and use
\[
 \partial_{z_j}=\tfrac12(\partial_{s_j}-i\partial_{t_j}),
 \qquad
 \partial_{\bar z_j}=\tfrac12(\partial_{s_j}+i\partial_{t_j}).
\]
The orientation is the complex orientation.  In one complex variable,
\(dA=ds\,dt=\frac{i}{2}dz\wedge d\bar z\).
There are no boundary conditions, boundary growth bounds, or support
restrictions on the fields.

The space \(C^\infty(U;\mathbb C)\) consists of complex-valued functions
with continuous real partial derivatives of every order.  If
\(\gamma\in\mathbb N^4\), then \(D^\gamma\) means the corresponding
product of derivatives in \((s_1,t_1,s_2,t_2)\), and
\(|\gamma|=\sum_j\gamma_j\).  For \(\alpha\in\mathbb N^2\), use
\(\partial_z^\alpha=\partial_{z_1}^{\alpha_1}
\partial_{z_2}^{\alpha_2}\),
\(\partial_{\bar z}^\alpha=\partial_{\bar z_1}^{\alpha_1}
\partial_{\bar z_2}^{\alpha_2}\),
\(|\alpha|=\alpha_1+\alpha_2\), and
\(\alpha!=\alpha_1!\alpha_2!\).

\begin{definition}[The Dolbeault algebra]
\label{def:tb-dolbeault}
Put \(dz_j=ds_j+i\,dt_j\) and \(d\bar z_j=ds_j-i\,dt_j\).
The generators \(d\bar z_1,d\bar z_2\) anticommute under exterior
multiplication, and a repeated generator gives zero.  For an increasing
set \(I=\{i_1<\cdots<i_q\}\subset\{1,2\}\), set
\(d\bar z_I=d\bar z_{i_1}\wedge\cdots\wedge d\bar z_{i_q}\),
with \(d\bar z_\varnothing=1\).  Define the Dolbeault carrier by
\[
 \mathcal A^q=\Omega^{0,q}(U)
   =\bigoplus_{\substack{I\subset\{1,2\}\\|I|=q}}
               C^\infty(U;\mathbb C)\,d\bar z_I
                         \qquad(0\le q\le2),
\]
and set \(\mathcal A^q=0\) otherwise.  Thus degree one consists of
\(a\,d\bar z_1+b\,d\bar z_2\), and degree two consists of
\(c\,d\bar z_1\wedge d\bar z_2\), with smooth coefficients.
The integer \(q\) is the cohomological degree.  An element in a single
degree is called homogeneous.  The differential and product are
\[
 d=\bar\partial=\sum_{j=1}^2d\bar z_j\wedge\partial_{\bar z_j},
 \qquad (a,b)\longmapsto a\wedge b.
\]
A differential graded algebra has a degree-one differential satisfying
\(d^2=0\) and
\(d(ab)=(da)b+(-1)^{|a|}a(db)\).
Here the commutation of the coefficient derivatives and the antisymmetry
of exterior multiplication give both identities.
\end{definition}

The holomorphic symplectic form is \(dz_1\wedge dz_2\).  Fix its bracket
normalization by
\[
 \{f,g\}=\partial_{z_1}f\,\partial_{z_2}g
          -\partial_{z_2}f\,\partial_{z_1}g.
\]
Thus the Hamiltonian derivation is
\(X_f=\{f,-\}=-(\partial_{z_2}f)\partial_{z_1}
 +(\partial_{z_1}f)\partial_{z_2}\), and
\(\iota_{X_f}(dz_1\wedge dz_2)=-\partial f\), where
\(\partial f=\sum_j(\partial_{z_j}f)\,dz_j\).
Here \(\iota_X\) inserts \(X\) into the first argument of a form:
\(\iota_X(dz_1\wedge dz_2)=(Xz_1)dz_2-(Xz_2)dz_1\).
The displayed sign follows directly from the two coefficients of \(X_f\).
Set
\begin{equation}
 [f\,d\bar z_I,g\,d\bar z_J]
       =\{f,g\}\,d\bar z_I\wedge d\bar z_J.
 \label{eq:tb-form-bracket}
\end{equation}

\begin{definition}[Differential graded Lie and Poisson algebras]
\label{def:tb-dglie}
A differential graded Lie algebra is a graded vector space with a
degree-one differential \(d\), \(d^2=0\), and a complex-bilinear bracket
of degree zero.  For homogeneous elements it satisfies
\[
 [a,b]=-(-1)^{|a||b|}[b,a],\qquad
 [a,[b,c]]=[[a,b],c]+(-1)^{|a||b|}[b,[a,c]],
\]
and
\[
 d[a,b]=[da,b]+(-1)^{|a|}[a,db].
\]
A differential graded Poisson algebra also has an associative product
satisfying
\[
 ab=(-1)^{|a||b|}ba,\qquad
 d(ab)=(da)b+(-1)^{|a|}a(db),\qquad
 [a,bc]=[a,b]c+(-1)^{|a||b|}b[a,c].
\]
Thus both the differential and the bracket act by graded derivations.
A morphism has degree zero, commutes with the differential,
and preserves the bracket and any specified product.
\end{definition}

\begin{lemma}[The coefficient bracket]
\label{lem:tb-poisson}
The product, differential, and bracket above make \(\mathcal A\) a unital
differential graded Poisson algebra.
\end{lemma}
\begin{proof}
For coefficient functions, the product rule proves the biderivation
identity.  In the cyclic Jacobi sum, expansion gives terms with one second
derivative and two first derivatives.  Each second-derivative term cancels
the term obtained by interchanging its two remaining factors.  For example,
the terms containing \(f_{z_1z_1}\) are
\(f_{z_1z_1}g_{z_2}h_{z_2}\) and its negative.  The mixed second
derivatives cancel because \(f_{z_1z_2}=f_{z_2z_1}\).
The same cancellations apply to the second derivatives of \(g\) and \(h\).
This proves the ordinary Jacobi identity.

Exterior multiplication of the constant forms supplies the Koszul signs in
the graded antisymmetry, Jacobi, and biderivation identities.
Since \(\partial_{\bar z_j}\) commutes with every \(\partial_{z_i}\),
the product rule for \(\bar\partial\) gives the differential bracket
identity.  The constant function one is a unit and is closed.
\end{proof}

Put
\[
 R=\mathbb C[[x,y]],\qquad \mathfrak m=(x,y),\qquad
 \mathfrak h=R/\mathbb C.
\]
An element of \(R\) is a coefficient family
\(F=\sum_{a,b\in\mathbb N}F_{ab}x^ay^b\).  Equality and addition are
coefficientwise.  Multiplication is finite coefficient convolution:
\[
 [x^ay^b](FG)=\sum_{i=0}^a\sum_{j=0}^bF_{ij}G_{a-i,b-j}.
\]
The notation \([x^ay^b]F\) means \(F_{ab}\).  For \(N\ge0\),
\[
 \mathfrak m^N=\left\{\sum_{a,b}F_{ab}x^ay^b:
                         F_{ab}=0\text{ when }a+b<N\right\},
 \qquad\mathfrak m^0=R.
\]
This agrees with the ideal power of \(\mathfrak m=(x,y)\).  Products
of \(N\) generators have degree at least \(N\).  Conversely, group
each term of degree at least \(N\) by the factor
\(x^{\min(a,N)}y^{N-\min(a,N)}\).  This writes every such series as
a sum of \(N+1\) monomials of degree \(N\) times formal series.
The symbols \(x,y\) are formal coordinates corresponding to \(z_1,z_2\).
Formal differentiation means
\(\partial_xF=\sum_{a\ge1,b\ge0}aF_{ab}x^{a-1}y^b\), and similarly
for \(\partial_y\).  These operations give the same Poisson formula
\(\{f,g\}=f_xg_y-f_yg_x\) on \(R\).
Represent each class in \(\mathfrak h\) by its unique series in
\(\mathfrak m\).  Its bracket is the formal bracket followed by removal
of the constant term.  In particular,
\begin{equation}
 \{x,y\}=0\text{ in }\mathfrak h,
 \qquad \{x,y^2\}=2y,
 \qquad \operatorname{ad}_x=\partial_y\pmod{\mathbb C}.
 \label{eq:tb-translations}
\end{equation}
Thus quotienting scalar Hamiltonians retains translations.

\begin{definition}[Pure holomorphic Taylor extraction]
\label{def:tb-taylor}
For \(\alpha=(\alpha_1,\alpha_2)\in\mathbb N^2\), write
\(\alpha!=\alpha_1!\alpha_2!\) and \(|\alpha|=\alpha_1+\alpha_2\).
Define
\begin{equation}
 Jf=\sum_{\alpha\in\mathbb N^2}
     \frac{\partial_z^\alpha f(0)}{\alpha!}
              x^{\alpha_1}y^{\alpha_2}
 \quad(f\in\mathcal A^0),
 \qquad J|_{\mathcal A^{q>0}}=0.
 \label{eq:tb-taylor}
\end{equation}
The target \(R\) has degree zero and zero differential.
\end{definition}

\begin{proposition}[Taylor extraction and scalar Hamiltonians]
\label{prop:tb-one}
The map \(J:\mathcal A\to R\) is a unital morphism of differential
graded Poisson algebras.  It induces a differential graded Lie morphism
\[
 \bar J:\mathcal L:=\mathcal A/(\mathbb C\cdot1)
                         \longrightarrow\mathfrak h.
\]
Neither quotient inherits the stated associative multiplication.
\end{proposition}
\begin{proof}
For every \(\alpha\), the product rule gives
\[
 \frac{\partial_z^\alpha(fg)(0)}{\alpha!}
 =\sum_{\beta+\gamma=\alpha}
   \frac{\partial_z^\beta f(0)}{\beta!}
   \frac{\partial_z^\gamma g(0)}{\gamma!}.
\]
Hence \(J(fg)=J(f)J(g)\).  Direct reindexing gives
\[
 J(\partial_{z_1}f)=\partial_xJf,
 \qquad J(\partial_{z_2}f)=\partial_yJf.
\]
These identities prove bracket compatibility on degree-zero inputs.
If either input has positive degree, both sides of product and bracket
compatibility vanish.  Also \(J\bar\partial=0\) and \(J1=1\).

The line \(\mathbb C\cdot1\) is a central Lie ideal and is preserved
by the differential.  The quotient Lie map therefore exists.
That line is not an associative ideal: a descended product would have
\([1][f]=0\), whereas the class of \(1f\) is \([f]\ne0\) for
nonconstant \(f\).  The same contradiction applies to \(R/\mathbb C\).
\end{proof}

The Lie quotients and the Poisson algebras fit into the commutative diagram
\[
\begin{tikzcd}[column sep=large]
 \mathcal A \arrow[r,"J"] \arrow[d] & R \arrow[d]\\
 \mathcal L \arrow[r,"\bar J"] & \mathfrak h .
\end{tikzcd}
\]
The vertical maps preserve the differential and Lie bracket.  An
associative quotient product on the bottom row is excluded by
Proposition~\ref{prop:tb-one}.

\section{Smooth realization and holomorphic convergence}
\label{sec:tb-realization}

The Taylor map records derivatives, without imposing bounds on their
growth.  A direct cutoff construction separates this property from
holomorphic convergence.

For a continuous function, its support is the closure of the set where
it is nonzero.  The notation \(C_c^\infty(U)\) denotes smooth functions
whose support is a compact subset of \(U\).

\begin{lemma}[Limits of smooth derivatives]
\label{lem:tb-smooth-limits}
Let \(W\subset\mathbb R^m\) be open.  Suppose \(f_N\in C^\infty(W)\)
and every sequence \(D^\gamma f_N\) converges uniformly on compact
subsets to a function \(g_\gamma\).  Then \(f=g_0\) is smooth and
\(D^\gamma f=g_\gamma\) for every \(\gamma\).
The same assertion holds for nets in place of sequences.
\end{lemma}
\begin{proof}
Each \(g_\gamma\) is continuous, since it is a locally uniform limit
of continuous functions.  Fix a closed coordinate box contained in
\(W\) and a real coordinate direction \(e_j\).  For a segment in
that box, the fundamental theorem of calculus gives
\[
 D^\gamma f_N(x+te_j)-D^\gamma f_N(x)
   =\int_0^tD^{\gamma+e_j}f_N(x+se_j)\,ds.
\]
Uniform convergence on the segment permits passage to the limit.  The
error of the integral is bounded by \(|t|\) times the uniform error of
its integrand.  Hence
\[
 g_\gamma(x+te_j)-g_\gamma(x)
                =\int_0^t g_{\gamma+e_j}(x+se_j)\,ds.
\]
Dividing by \(t\) and using continuity of the integrand proves
\(\partial_jg_\gamma=g_{\gamma+e_j}\).  Starting with \(\gamma=0\)
and iterating proves the assertion.  Every point lies in such a box.
The proof uses eventual uniform bounds only, so it also applies to nets.
\end{proof}

\begin{proposition}[Realization of arbitrary formal jets]
\label{prop:tb-two}
For every \(F=\sum_\alpha a_\alpha x^{\alpha_1}y^{\alpha_2}\in R\),
there is \(f\in C_c^\infty(U)\) such that
\begin{equation}
 \partial_z^\alpha\partial_{\bar z}^{\beta}f(0)
 =\begin{cases}
   \alpha!a_\alpha,&\beta=0,\\
   0,&\beta\ne0.
  \end{cases}
 \label{eq:tb-exact-jets}
\end{equation}
In particular, \(J\) and \(\bar J\) are surjective.
\end{proposition}
\begin{proof}
Choose \(0<\delta<1\) with \(\overline{B(0,\delta)}\subset U\).
An explicit cutoff on \(\mathbb R^4\) is
\[
 \eta(t)=\begin{cases}e^{-1/t},&t>0,\\0,&t\le0,\end{cases}
 \qquad
 \chi(w)=\frac{\eta(1-|w|^2)}
 {\eta(1-|w|^2)+\eta(|w|^2-\tfrac14)}.
\]
Every right derivative of \(\eta\) is a polynomial in \(1/t\) times
\(e^{-1/t}\), and tends to zero at zero.  Thus \(\eta\) and \(\chi\)
are smooth.  The denominator is positive everywhere.
The function \(\chi\) equals one on \(|w|\le1/2\), and vanishes
on \(|w|\ge1\).

Group the coefficients by homogeneous degree:
\[
 P_n(z)=\sum_{|\alpha|=n}a_\alpha z^\alpha.
\]
For real multi-indices \(\gamma\), set
\[
 C_{n,k}=\max_{|\gamma|\le k}
       \sup_{w\in\mathbb R^4}|D^\gamma(\chi P_n)(w)|.
\]
These numbers are finite.  For \(n\ge1\), choose
\[
 0<\varepsilon_n\le
 \min\left\{\delta,2^{-n},
       \min_{0\le k<n}
       \left(\frac{2^{-n}}{1+C_{n,k}}\right)^{1/(n-k)}\right\}.
\]
Define
\[
 f_0(z)=a_0\chi(z/\delta),\qquad
 f_n(z)=\chi(z/\varepsilon_n)P_n(z)\quad(n\ge1).
\]
Real homogeneity gives
\[
 D^\gamma f_n(z)=\varepsilon_n^{n-|\gamma|}
            D^\gamma(\chi P_n)(z/\varepsilon_n).
\]
Since \(\varepsilon_n<1\), for \(|\gamma|\le k<n\) one has
\(\varepsilon_n^{n-|\gamma|}\le\varepsilon_n^{n-k}\).  Therefore
\begin{equation}
 \max_{|\gamma|\le k}\|D^\gamma f_n\|_\infty\le2^{-n}
                  \qquad(k<n).
 \label{eq:tb-cutoff-estimate}
\end{equation}

For every fixed \(k\), the tail \(n>k\) converges uniformly with all
real derivatives of order at most \(k\).  The finitely many remaining
terms are smooth.  Lemma~\ref{lem:tb-smooth-limits}, applied to the
partial sums, shows that \(f=\sum_{n\ge0}f_n\) is smooth with the
termwise real derivatives.  All supports
lie in \(\overline{B(0,\delta)}\), so \(f\in C_c^\infty(U)\).

Each individual \(f_n\) equals \(P_n\) on a neighborhood of zero.
Uniform derivative convergence permits termwise evaluation there.
A derivative of total order \(k\) receives a contribution only from
\(P_k\).  All antiholomorphic derivatives of this polynomial vanish.
The remaining holomorphic derivative is exactly \(\alpha!a_\alpha\).
This proves~\eqref{eq:tb-exact-jets}.
\end{proof}

The constructed \(\bar\partial f\) is flat at the origin.  Flatness
means that every real derivative of each coefficient vanishes there.
It does not imply that \(\bar\partial f\) vanishes on a neighborhood.
The choices of radii depend on the prescribed coefficients; no linear
extension operator is asserted.

\subsection{Cauchy integration and holomorphicity}
\label{sec:tb-cauchy}

\begin{definition}[Holomorphic functions]
\label{def:tb-holomorphic}
A function on an open subset of \(\mathbb C^m\), with \(m=1\) or
\(2\), is holomorphic if near every point \(p\) it is an absolutely
convergent complex power series \(\sum_\alpha a_\alpha(z-p)^\alpha\).
Absolute convergence is required on some polydisc centered at \(p\).
The space of such functions on \(U\) is denoted \(\mathcal O(U)\).
The equivalence with smooth solutions of the Cauchy--Riemann equations
is proved below.
\end{definition}

For a piecewise smooth oriented curve \(\gamma:[a,b]\to\mathbb C\),
the contour integral means
\(\int_\gamma f(\zeta)\,d\zeta
 =\int_a^bf(\gamma(t))\gamma'(t)\,dt\).
A positively oriented circle is traversed counterclockwise.

\begin{lemma}[Cauchy--Pompeiu formula]
\label{lem:tb-cauchy-pompeiu}
Let \(f\) be smooth on a neighborhood of the closed disc
\(\overline{D(0,r)}\), and let \(|w|<r\).  Then
\begin{equation}
 f(w)=\frac1{2\pi i}\int_{|\zeta|=r}
             \frac{f(\zeta)}{\zeta-w}\,d\zeta
       -\frac1\pi\int_{|\zeta|<r}
             \frac{\partial_{\bar\zeta}f(\zeta)}{\zeta-w}\,dA(\zeta).
 \label{eq:tb-cauchy-pompeiu}
\end{equation}
In particular, every \(\psi\in C_c^\infty(\mathbb C)\) satisfies
\begin{equation}
 -\int_{\mathbb C}\frac{\partial_{\bar\zeta}\psi(\zeta)}
                             {\pi\zeta}\,dA(\zeta)=\psi(0).
 \label{eq:tb-test-identity}
\end{equation}
\end{lemma}
\begin{proof}
For a smooth function \(g\), direct calculation in \(\zeta=s+it\)
gives
\[
 d(g\,d\zeta)=(i\partial_sg-\partial_tg)\,ds\wedge dt
                         =2i\,\partial_{\bar\zeta}g\,dA.
\]
The planar integration identity associated to this equality follows
from one-variable real calculus.  Integrate \(\partial_tP\) on each
vertical interval of a region, and integrate \(\partial_sQ\) on each
horizontal interval.  The fundamental theorem of calculus gives the
endpoint values.  Integration in the remaining variable yields
\[
 \int_{\partial A}(P\,ds+Q\,dt)
                 =\int_A(\partial_sQ-\partial_tP)\,ds\,dt.
\]
For a disc with an interior disc removed, these slices have at most
two component intervals.  Their endpoint terms are exactly the positively
oriented outer circle and the negatively oriented inner circle.  Splitting
at the finitely many extrema of the circles justifies the graph
parameterizations; taking endpoint limits gives the displayed formula.
All derivatives are continuous on the compact punctured region, so the
iterated integrations are valid.

Apply this identity on
\(A_\epsilon=D(0,r)\setminus\overline{D(w,\epsilon)}\) to
\(g(\zeta)=f(\zeta)/(\zeta-w)\), where
\(0<\epsilon<r-|w|\).  Away from \(w\), direct differentiation gives
\(\partial_{\bar\zeta}(1/(\zeta-w))=0\).  Consequently
\[
 \begin{split}
 \frac1{2\pi i}\int_{|\zeta|=r}\frac{f(\zeta)}{\zeta-w}\,d\zeta
 &-\frac1{2\pi i}\int_{|\zeta-w|=\epsilon}
                         \frac{f(\zeta)}{\zeta-w}\,d\zeta\\
 &=\frac1\pi\int_{A_\epsilon}
                \frac{\partial_{\bar\zeta}f(\zeta)}{\zeta-w}\,dA.
 \end{split}
\]
Both written circles are counterclockwise; the minus sign accounts for
the inner boundary orientation.  Parameterizing the small circle gives
\[
 \frac1{2\pi i}\int_{|\zeta-w|=\epsilon}
            \frac{f(\zeta)}{\zeta-w}\,d\zeta
       =\frac1{2\pi}\int_0^{2\pi}f(w+\epsilon e^{it})\,dt
       \longrightarrow f(w).
\]
The omitted area integral tends to zero.  If
\(|\partial_{\bar\zeta}f|\le M\) there, its absolute value is at
most \(2M\epsilon\), since
\(\int_{|v|<\epsilon}|v|^{-1}dA(v)=2\pi\epsilon\).
Letting \(\epsilon\to0\) proves~\eqref{eq:tb-cauchy-pompeiu}.

For compactly supported \(\psi\), choose \(r\) so large that its
support lies inside the circle and put \(w=0\).  The boundary term
vanishes, giving~\eqref{eq:tb-test-identity}.
\end{proof}

\begin{lemma}[Cauchy--Riemann equations and convergent expansions]
\label{lem:tb-cr-expansion}
A function on an open subset of \(\mathbb C\) or \(\mathbb C^2\)
is holomorphic exactly when it is smooth and all its antiholomorphic
partial derivatives vanish.  If \(f\) is smooth on a neighborhood of
a closed bidisc and satisfies these equations, then
\begin{equation}
 f(w_1,w_2)=\frac1{(2\pi i)^2}
  \int_{|\zeta_1|=r_1}\int_{|\zeta_2|=r_2}
  \frac{f(\zeta_1,\zeta_2)}
             {(\zeta_1-w_1)(\zeta_2-w_2)}\,d\zeta_2\,d\zeta_1
 \label{eq:tb-bidisc-cauchy}
\end{equation}
for \(|w_i|<r_i\).  Its expansion at zero has coefficients
\begin{equation}
 \begin{aligned}
 a_\alpha&=\frac1{(2\pi i)^2}
 \int_{|\zeta_1|=r_1}\int_{|\zeta_2|=r_2}
 \frac{f(\zeta_1,\zeta_2)}
      {\zeta_1^{\alpha_1+1}\zeta_2^{\alpha_2+1}}\,d\zeta_2\,d\zeta_1,\\
 |a_\alpha|&\le M r_1^{-\alpha_1}r_2^{-\alpha_2}.
 \end{aligned}
 \label{eq:tb-cauchy-coefficients}
\end{equation}
where \(M=\sup_{|\zeta_i|=r_i}|f|\).  These coefficients equal
\(\partial_z^\alpha f(0)/\alpha!\).
\end{lemma}
\begin{proof}
First consider one variable with \(\partial_{\bar z}f=0\).
The area term in~\eqref{eq:tb-cauchy-pompeiu} is zero.  For
\(0<s<r\) and \(|w|\le s\), the geometric expansion
\[
 \frac1{\zeta-w}=\sum_{k\ge0}\frac{w^k}{\zeta^{k+1}}
                         \qquad(|\zeta|=r)
\]
converges uniformly.  Its remainder after degree \(N\) has absolute
value at most \((s/r)^{N+1}/(r-s)\).  Termwise contour integration
therefore gives
\[
 f(w)=\sum_{k\ge0}a_kw^k,\qquad
 a_k=\frac1{2\pi i}\int_{|\zeta|=r}
                    \frac{f(\zeta)}{\zeta^{k+1}}\,d\zeta,
 \qquad |a_k|\le Mr^{-k}.
\]
Every fixed real derivative of order \(j\) introduces at most a constant
times \(k^j\).  Its terms on \(|w|\le s\) are bounded by a constant
times \(k^j(s/r)^k\).  This series converges: the ratio of consecutive
terms tends to \(s/r<1\).  Lemma~\ref{lem:tb-smooth-limits} justifies
termwise differentiation of all orders.  Evaluation at zero gives
\(a_k=\partial_z^k f(0)/k!\).

In two variables, apply the one-variable boundary formula first in
\(z_1\).  Apply it again in \(z_2\) to the boundary coefficients.
The two finite contour integrals yield~\eqref{eq:tb-bidisc-cauchy}.
Expanding both kernels uniformly on smaller bidiscs gives
\eqref{eq:tb-cauchy-coefficients}.  Products of the preceding geometric
derivative estimates prove convergence with every real derivative and
the stated Taylor coefficient identity.  Translating the center proves
the local expansion at every point of an open set.

Conversely, suppose a function is given locally by an absolutely
convergent power series.  Choose positive radii \(s_i\) strictly inside
its convergence polydisc.  Absolute convergence at these radii means
\(\sum_\alpha|a_\alpha|s^\alpha<\infty\).
For any smaller radii \(t_i<s_i\), each differentiated term is bounded
by \(|a_\alpha|s^\alpha\) times a bounded polynomial-geometric factor
in \(\alpha\).  The differentiated series therefore converge uniformly
on the smaller closed polydisc.  Lemma~\ref{lem:tb-smooth-limits}
shows that the sum is smooth.  Each monomial has zero antiholomorphic
derivatives, so the sum satisfies the Cauchy--Riemann equations.
\end{proof}

\begin{definition}[Series converging on the specified polydisc]
\label{def:tb-convergent}
Let \(R_U\subset R\) consist of the series whose coefficients satisfy
\begin{equation}
 \sum_{\alpha\in\mathbb N^2}|a_\alpha|
             r_1^{\alpha_1}r_2^{\alpha_2}<\infty
 \quad\text{for every }0<r_i<R_i.
 \label{eq:tb-convergence}
\end{equation}
The domain \(U\) remains fixed in this definition.
\end{definition}

\begin{proposition}[The holomorphic image]
\label{prop:tb-three}
The map \(J:\mathcal O(U)\to R_U\) is bijective.
The subspace \(R_U\) is proper in \(R\).
For holomorphic germs at zero, the corresponding image is
\(\mathbb C\{x,y\}\), the series converging on some neighborhood of zero.
\end{proposition}
\begin{proof}
Choose radii \(r_i<s_i<R_i\).  Applying
Lemma~\ref{lem:tb-cr-expansion} on the closed bidisc of radii \(s_i\)
gives
\[
 a_\alpha=\frac1{(2\pi i)^2}
  \int_{|\zeta_1|=s_1}\int_{|\zeta_2|=s_2}
  \frac{f(\zeta_1,\zeta_2)}
       {\zeta_1^{\alpha_1+1}\zeta_2^{\alpha_2+1}}
                         \,d\zeta_2\,d\zeta_1.
\]
If \(M_s\) is the maximum of \(|f|\) on this product of circles, then
\[
 |a_\alpha|\le M_s s_1^{-\alpha_1}s_2^{-\alpha_2}.
\]
Summing the geometric series proves~\eqref{eq:tb-convergence}.
Expansion of both Cauchy kernels proves that the series equals \(f\)
on the smaller polydisc.  These polydiscs exhaust \(U\), proving
injectivity and the asserted image inclusion.  The circles have the
positive orientations of~\eqref{eq:tb-bidisc-cauchy}.

Conversely, condition~\eqref{eq:tb-convergence} gives absolute uniform
convergence on each smaller closed polydisc.  Its derivatives also converge
there.  Indeed, choose larger radii \(s_i<R_i\).  Every derivative
introduces only polynomial factors in \(\alpha_i\), and these factors
are bounded after multiplication by \((r_i/s_i)^{\alpha_i}\).
Lemma~\ref{lem:tb-smooth-limits} justifies termwise differentiation.
The antiholomorphic derivatives vanish, so
Lemma~\ref{lem:tb-cr-expansion} gives a holomorphic function with the
prescribed coefficients.

For \(F=\sum_{n\ge1}n!x^n\) and every \(r>0\), the ratio of
successive absolute terms is \((n+1)r\).  The terms do not tend to zero.
Hence \(F\notin R_U\), and it is not the Taylor series of a holomorphic
germ.  Applying the preceding argument on varying sufficiently small
polydiscs proves the assertion about germs.
\end{proof}

The quantifier in~\eqref{eq:tb-convergence} retains the specified domain.
Replacing \(\mathcal O(U)\) by germs changes the image and does not
remove the factorial counterexample.  The differential must detect this
failure of holomorphic realization.

\section{The Dolbeault complex and its Taylor kernel}
\label{sec:tb-cohomology}

The local equation \(\partial_{\bar z}u=a\) admits an integral solution
when \(a\) has compact support.  To solve it on an open disc with arbitrary
boundary growth, the far-annulus contributions require holomorphic
corrections.  This construction also controls smooth parameters.

\begin{lemma}[A one-variable primitive with parameters]
\label{lem:tb-parametric}
Let \(W\) be an open disc and let \(a\in C^\infty(D(0,R)\times W)\).
There is \(u\in C^\infty(D(0,R)\times W)\) with
\(\partial_{\bar z}u=a\).  If \(a\) is holomorphic in the parameter
\(w\), then \(u\) can be chosen holomorphic in \(w\).
\end{lemma}
\begin{proof}
For \(\phi\in C_c^\infty(\mathbb C)\), define
\begin{equation}
 T\phi(z)=\frac1\pi\int_{\mathbb C}
                 \frac{\phi(\zeta)}{z-\zeta}\,dA(\zeta).
 \label{eq:tb-cauchy-transform}
\end{equation}
The kernel is locally integrable, as the polar-coordinate calculation
in Lemma~\ref{lem:tb-cauchy-pompeiu} shows.  To prove smoothness, write
\[
 T\phi(z)=\int_{\mathbb C}\frac{\phi(z-v)}{\pi v}\,dA(v),
\]
and place every derivative on \(\phi\).  For \(z\) in a compact set,
only a bounded set of \(v\) contributes.  The difference quotients of
each derivative of \(\phi\) converge uniformly there, by its continuous
next derivatives and the one-variable fundamental theorem of calculus.
The error of their integrals is at most their uniform error times the
integral of \(1/(\pi|v|)\) over that bounded set.  This proves
differentiation under the integral, to every real order.
It also proves smooth dependence on additional parameters when the
support in the solved variable lies in one fixed compact set.

Apply~\eqref{eq:tb-test-identity} to
\(\psi(v)=\phi(z-v)\).  The chain rule gives the pointwise identity
\[
 \partial_{\bar z}T\phi(z)
 =-\int_{\mathbb C}
       \frac{\partial_{\bar v}[\phi(z-v)]}{\pi v}\,dA(v)
 =\phi(z).
\]
If \(\phi\) depends holomorphically on a parameter \(w\),
differentiation in \(\bar w\) under the integral gives zero.
Lemma~\ref{lem:tb-cr-expansion} then
proves preservation of holomorphic parameter dependence.

Choose \(0<r_0<r_1<\cdots\uparrow R\) and smooth radial cutoffs
\(\rho_n\) equal to one on \(|z|\le r_n\), with support in
\(|z|<r_{n+1}\).  Put
\[
 \chi_0=\rho_0,\qquad \chi_n=\rho_n-\rho_{n-1}\quad(n\ge1).
\]
Then \(\sum_n\chi_n=1\) locally finitely.  Each \(\chi_na\) has
compact support in the solved variable, so
\(v_n=T_z(\chi_na)\) is defined.  For \(n\ge1\), \(v_n\) is
holomorphic on \(|z|<r_{n-1}\), because \(\chi_n\) vanishes there
and Lemma~\ref{lem:tb-cr-expansion} applies.

Choose an increasing compact exhaustion \(K_n\) of \(W\).
For \(n\ge3\), take a Taylor polynomial \(P_n(z,w)\) of \(v_n\)
in \(z\) so that
\begin{equation}
 \sup_{\substack{|z|\le r_{n-2}\\w\in K_n}}
     |D^\gamma(v_n-P_n)(z,w)|\le2^{-n}
                          \qquad(|\gamma|\le n).
 \label{eq:tb-annular-estimate}
\end{equation}
Here \(\gamma\) ranges over real derivatives in both variables.
To justify the choice, select a circle of radius strictly between
\(r_{n-2}\) and \(r_{n-1}\).  The coefficient estimates proved in
Lemma~\ref{lem:tb-cr-expansion} on
that circle are uniform over \(K_n\) for each of the finitely many
parameter derivatives.  Their geometric tails tend to zero, also after
the finitely many \(z\)-derivatives.  One polynomial degree therefore
gives all estimates in~\eqref{eq:tb-annular-estimate}.
Its coefficients are smooth in \(w\), and holomorphic when \(a\) is.
For the finitely many initial terms, set \(P_n=0\).

Every fixed compact subset of \(D(0,R)\times W\) eventually lies in
the regions of~\eqref{eq:tb-annular-estimate}.  For every derivative
order, the corresponding tail is bounded by \(\sum_n2^{-n}\).
Consequently
\[
 u=\sum_{n\ge0}(v_n-P_n)
\]
converges with all real derivatives on compact subsets.  Its smoothness
and differentiated sums follow from Lemma~\ref{lem:tb-smooth-limits}.
Since each
\(P_n\) is holomorphic in \(z\),
\[
 \partial_{\bar z}u=\sum_n\chi_na=a.
\]
The convergence preserves holomorphic parameter dependence.
No estimate at the boundary of the open discs was used.
\end{proof}

\begin{definition}[Cohomology and comparison maps]
\label{def:tb-cohomology}
For a complex \((C^\bullet,d)\) with \(d:C^q\to C^{q+1}\) and
\(d^2=0\), define
\[
 H^q(C)=\ker(d:C^q\to C^{q+1})/
                       \operatorname{im}(d:C^{q-1}\to C^q).
\]
Elements in the kernel are cocycles.  Elements in the image are
coboundaries.  A primitive of a degree-\(q\) cocycle \(\alpha\) is
an element \(\beta\in C^{q-1}\) with \(d\beta=\alpha\).
A chain map \(T:C\to C'\) is a degree-zero linear map satisfying
\(Td=d'T\).  It induces maps on cohomology by sending the class of a
cocycle to the class of its image.  It is a quasi-isomorphism when
these maps are isomorphisms in every degree.  A section of \(T\) is
a map \(S\) with \(TS=\mathrm{id}\); a chain-map section also
preserves degrees, is linear, and commutes with the differentials.
\end{definition}

\begin{proposition}[Global polydisc exactness]
\label{prop:tb-four}
As cohomology groups of complexes of complex vector spaces,
\[
 H^q(\mathcal A)=
 \begin{cases}\mathcal O(U),&q=0,\\0,&q>0,\end{cases}
 \qquad
 H^q(\mathcal L)=
 \begin{cases}\mathcal O(U)/\mathbb C,&q=0,\\0,&q>0.\end{cases}
\]
\end{proposition}
\begin{proof}
For a closed \((0,1)\)-form
\[
 \alpha=a\,d\bar z_1+b\,d\bar z_2,
 \qquad \partial_{\bar z_1}b=\partial_{\bar z_2}a,
\]
Lemma~\ref{lem:tb-parametric} gives \(u\) with
\(\partial_{\bar z_1}u=a\).  The function
\(b'=b-\partial_{\bar z_2}u\) satisfies
\(\partial_{\bar z_1}b'=0\).  Apply the lemma in \(z_2\), preserving
holomorphic dependence on \(z_1\), to obtain
\(\partial_{\bar z_2}v=b'\) and \(\partial_{\bar z_1}v=0\).
Then \(\bar\partial(u+v)=\alpha\).

For \(\alpha=c\,d\bar z_1\wedge d\bar z_2\), solve
\(\partial_{\bar z_1}u=c\).  The form \(u\,d\bar z_2\) satisfies
\[
 \bar\partial(u\,d\bar z_2)
       =c\,d\bar z_1\wedge d\bar z_2.
\]
Higher form degrees vanish by dimension.
Finally, a smooth function with \(\bar\partial f=0\) satisfies the
Cauchy--Riemann equations in each variable.
Lemma~\ref{lem:tb-cr-expansion} identifies these functions exactly with
\(\mathcal O(U)\).  Thus this is the degree-zero kernel.

Removing scalar constants quotients the degree-zero kernel by that line.
It does not change the image of \(\bar\partial\) in degree one, so
the positive-degree cohomology remains zero.
\end{proof}

\begin{corollary}[The Taylor defect]
\label{prop:tb-five}
The map on degree-zero cohomology is the proper injection
\[
 H^0(\bar J):\mathcal O(U)/\mathbb C\longrightarrow\mathfrak h,
 \qquad \operatorname{im}H^0(\bar J)=R_U/\mathbb C.
\]
In particular, \(\bar J\) is not a quasi-isomorphism.
For \(K=\ker(J:\mathcal A\to R)\),
\begin{equation}
 H^0(K)=0,\qquad H^1(K)\cong R/R_U,
 \qquad H^{q\ge2}(K)=0.
 \label{eq:tb-kernel-cohomology}
\end{equation}
The kernel of \(\bar J\) is canonically isomorphic to \(K\).
\end{corollary}
\begin{proof}
A degree-zero cocycle in \(K\) is holomorphic and has zero Taylor
series.  Proposition~\ref{prop:tb-three} makes it zero.
For \(F\in R\), choose a smooth lift \(f\) from
Proposition~\ref{prop:tb-two}, and define
\begin{equation}
 \delta(F)=[\bar\partial f]\in H^1(K).
 \label{eq:tb-connecting}
\end{equation}
Two lifts differ by an element of \(K^0\), so the class is independent
of the lift.  This also proves linearity, without a linear choice of lifts.
The equality \(\bar\partial f=\bar\partial k\) for some \(k\in K^0\)
holds exactly when \(f-k\) is holomorphic with Taylor series \(F\).
Thus \(\ker\delta=R_U\).

Every closed element of \(K^1=\mathcal A^1\) has a smooth primitive by
Proposition~\ref{prop:tb-four}.  Hence \(\delta\) is surjective.
In degree at least two, the same proposition gives a primitive of positive
degree, which automatically lies in \(K\).  This proves
\eqref{eq:tb-kernel-cohomology}.

If \([f]\in\ker\bar J\), then \(Jf\) is constant.  Replacing \(f\)
by \(f-f(0)\) identifies this degree-zero kernel with \(K^0\).
Positive degrees are unchanged, and the identification commutes with
the differential.
\end{proof}

All quotients in~\eqref{eq:tb-kernel-cohomology} are vector-space
quotients.  No Hausdorff quotient topology or ordinary Lie bracket on
\(R/R_U\) is asserted.  For the factorial series, any smooth lift
gives a nonzero class \([\bar\partial f]\) in \(H^1(K)\).
The relevant exact sequence is therefore
\[
 0\longrightarrow\mathcal O(U)\xrightarrow{J}R
       \xrightarrow{\delta}H^1(K)\longrightarrow0.
\]

\begin{corollary}[No section as a complex]
\label{cor:tb-no-chain-section}
Neither \(J\) nor \(\bar J\) has a chain-map section, even as a map
of complexes without topology.
\end{corollary}
\begin{proof}
Every formal series is a degree-zero cocycle.  A chain-map section would
lift it to a degree-zero cocycle in the Dolbeault complex, hence to a
holomorphic function.  The factorial series contradicts this conclusion.
For the quotient complex, closed degree-zero representatives are still
holomorphic, so the same contradiction applies.
\end{proof}

\section{Two formal topologies}
\label{sec:tb-topologies}

The algebraic map exists before any topology is chosen.  Continuity
depends on whether a small coefficient may vary or must vanish exactly.

\begin{definition}[Seminorm topologies]
\label{def:tb-seminorms}
A seminorm on a complex vector space \(E\) is a function
\(p:E\to[0,\infty)\) satisfying
\[
 p(ax)=|a|p(x),\qquad p(x+y)\le p(x)+p(y).
\]
A family of seminorms gives a neighborhood base at zero by the sets
\(\bigcap_{j=1}^k\{x:p_j(x)<\epsilon_j\}\), where \(k\) is finite
and \(\epsilon_j>0\).  Their translates define the topology.
A locally convex complex space is a vector space with such a topology.
It is Hausdorff if the seminorms separate points.  A net \((x_a)\) is
a family indexed by a directed set: every two indices have a common
upper bound.  It is Cauchy if, for every finite family and every positive tolerance, its
sufficiently late differences lie in the corresponding zero-neighborhood.
The space is complete if every Cauchy net converges in it.
A Fr\'echet space is a Hausdorff complete locally convex space whose
topology is generated by a countable family of seminorms.
\end{definition}

For a countable separating family \((p_j)_{j\ge1}\), the formula
\[
 d_E(x,y)=\sum_{j\ge1}2^{-j}\min\{1,p_j(x-y)\}
\]
is a metric giving the same topology.  The triangle inequality follows
from the seminorm inequalities and
\(\min(1,a+b)\le\min(1,a)+\min(1,b)\).
Controlling finitely many seminorms controls this sum up to its geometric
tail; making the sum small controls each specified seminorm.  This proves
the topology assertion and explains the metrizability of the spaces used
below.

\begin{definition}[Smooth, coefficient, and adic topologies]
\label{def:tb-topologies}
On smooth coefficient functions use seminorms
\[
 p_{K,r}(f)=\max_{|\gamma|\le r}\sup_{z\in K}|D^\gamma f(z)|,
\]
where \(K\Subset U\), \(r\ge0\), and \(\gamma\) is real.
Use the finite product topology for the coefficients of a differential
form.  A countable compact exhaustion makes each \(\mathcal A^q\)
a metrizable Fr\'echet space.

The usual coefficient topology on \(R\) is
\(\prod_{\alpha\in\mathbb N^2}\mathbb C\), with usual complex factors.
The adic topology has the subspaces \(\mathfrak m^N\) as a neighborhood
base at zero and is taken over \(\mathbb C_{\mathrm{disc}}\).
Give \(\mathfrak h\) the corresponding quotient topology, represented
on \(\mathfrak m\).
\end{definition}

Smooth spaces are complete because a Cauchy net converges uniformly on
compact sets in every derivative.  For each compact set, completeness
of the scalar field gives the pointwise limit, and the uniform Cauchy
bound passes to a uniform limit bound.  The limits are compatible under
restriction.  Lemma~\ref{lem:tb-smooth-limits} makes them the derivatives
of one smooth function.  The coefficient product is complete
because its scalar factors are complete.  The adic space is complete
because a compatible family of finite coefficient assignments gives one
formal series.  These are different completeness assertions.

The scalar constants form a closed smooth subspace.  The normalization
\begin{equation}
 [f]\longmapsto f-f(0)
 \label{eq:tb-normalization}
\end{equation}
identifies \(\mathcal L^0\) continuously with the closed subspace of
functions vanishing at zero.  Thus the finite graded sum underlying
\(\mathcal L\) is also metrizable and complete.

\begin{proposition}[Continuity of the Taylor map]
\label{prop:tb-six}
The maps \(J\) and \(\bar J\) are continuous for the usual coefficient
topology.  They are discontinuous into the adic target when the source
retains its smooth topology.  The brackets are jointly continuous on the
smooth source and for both formal target topologies.
\end{proposition}
\begin{proof}
Every Taylor coefficient is a derivative evaluation at zero, which is a
continuous smooth linear functional.  A map to a product is continuous
exactly when its coordinate maps are continuous.  Normalization
\eqref{eq:tb-normalization} proves the quotient assertion.

Take nonzero usual complex numbers \(t_n\to0\).  Then
\(t_nz_1\to0\) in every smooth seminorm, whereas
\(J(t_nz_1)=t_nx\notin\mathfrak m^2\).  The formal images do not
converge to zero adically.  This example also applies modulo constants.

An output coefficient of a formal bracket is a finite polynomial in input
coefficients, proving usual joint continuity.  More precisely, its output
through degree \(N\) depends only on inputs through degree \(N+1\).
For the adic topology this follows from
\begin{equation}
 \{\mathfrak m^a,\mathfrak m^b\}
       \subset\mathfrak m^{a+b-2},\qquad a,b\ge1,
 \label{eq:tb-order-estimate}
\end{equation}
interpreted modulo constants when necessary.  Perturbing both inputs by
series of order at least \(N+2\) leaves their bracket unchanged through
degree \(N\), including the mixed perturbation terms.

On the smooth source, the product rule gives
\[
 p_{K,r}(\{f,g\})\le C_r p_{K,r+1}(f)p_{K,r+1}(g).
\]
The same bound on finitely many form coefficients proves joint continuity
of~\eqref{eq:tb-form-bracket}.  It also proves continuity of the exterior
product; differentiation is continuous by the defining seminorms.
\end{proof}

The identity from the adic formal space to the usual coefficient space is
continuous: the inverse image of a product neighborhood contains a tail.
The reverse identity is discontinuous.  For example, the condition that
the linear coefficient vanish is adically open but is not usually open.
Scalar multiplication \(\mathbb C\times R_{\mathrm{adic}}\to
R_{\mathrm{adic}}\) is not continuous for usual \(\mathbb C\), as
the same sequence \(t_nx\) proves.  The discrete coefficient field in
Definition~\ref{def:tb-topologies} is therefore part of the category.

Polynomials realize every finite coefficient assignment, so they are
dense in both formal topologies.  This proves that the proper subspace
\(R_U\) is dense in both.  Density does not give a holomorphic realization
of a divergent series.

\begin{proposition}[No continuous linear section]
\label{prop:tb-no-linear-section}
For the usual coefficient topology, neither \(J\) nor \(\bar J\)
has a continuous complex-linear section into its smooth source.
\end{proposition}
\begin{proof}
Suppose \(S:R\to C^\infty(U)\) were such a section.
Fix a compact ball \(K\) containing zero in its interior.  Continuity
for \(p_{K,0}\) gives a coefficient neighborhood \(V\) on which
\(p_{K,0}(SF)<1\).  The neighborhood \(V\) contains an entire tail
subspace \(\mathfrak m^N\).  Scaling any vector in this tail shows
that \(SF\) vanishes on \(K\).  Its Taylor series is then zero,
contradicting \(JSF=F\) for a nonzero tail vector.
For \(\bar J\), compose the proposed section with normalization
\eqref{eq:tb-normalization} and repeat the argument on \(\mathfrak h\).
\end{proof}

The smooth surjection supplies individual lifts.  Its failure to split as
a complex and its failure to split continuously as a linear map are
distinct obstructions.  Neither obstructs precomposition of scalar forms
that use only finitely many Taylor coefficients.

\section{Continuous scalar cochains and their completion}
\label{sec:tb-cochains}

A scalar cochain evaluates a finite tuple of symmetry parameters.
Joint continuity forces such an evaluation to depend on only finitely
many coefficients, even though the algebra contains all translations.

\begin{proposition}[Finite jet support]
\label{prop:tb-seven}
Let \(n\ge1\).  Every jointly continuous complex-multilinear scalar
form \(\omega:\mathfrak h^n\to\mathbb C\) has an integer bound
\(N\ge0\) such that
\begin{equation}
 \omega(f_1,\ldots,f_n)=0
 \quad\text{if any }f_i\in\mathfrak m^{N+1}.
 \label{eq:tb-finite-support}
\end{equation}
This holds for usual coefficients with usual scalar values, and for adic
coefficients with either discrete or usual scalar values.  Conversely,
every form satisfying~\eqref{eq:tb-finite-support} is continuous in these
cases.  Arity zero consists of the scalar constants.
\end{proposition}
\begin{proof}
The converse follows by factoring through finite coefficient spaces.
For the forward direction, continuity at the origin gives neighborhoods
\(U_i\) on whose product \(|\omega|<1\); use \(\{0\}\) instead for
discrete scalar values.  Every \(U_i\) contains a tail subspace.
Choose a common tail \(T=\mathfrak m^M\).
Since every scalar multiple of an element of \(T\) is still in \(T\),
multilinearity implies \(\omega|_{T^n}=0\).  Otherwise scaling one
argument would produce arbitrarily large values in that product.

Proceed by induction on \(n\).  The preceding argument proves the case
\(n=1\).  Split \(\mathfrak h=H\oplus T\), with \(H\) the finite
space of the preceding homogeneous degrees.  Expanding each argument
into head and tail parts gives finitely many terms.  The all-tail term
vanishes.  In every other term, expand the head arguments in a finite
basis of \(H\).  Fixing these basis vectors leaves a continuous form
of fewer arguments.  By induction it has finite jet support.
Taking the maximum of the finitely many resulting bounds proves
\eqref{eq:tb-finite-support}.
\end{proof}

Index the nonconstant monomials by
\(\mathcal I=\{(a,b)\in\mathbb N^2:a+b\ge1\}\), and put
\(e_I=x^ay^b\).  Define the coefficient functional by
\(e_I^\vee(F)=[x^ay^b]F\).  These monomials are a coordinate family
for the formal product, not a Hamel basis: a formal series can have
infinitely many nonzero coordinates.  The continuous scalar dual is,
as a vector space,
\[
 D=\bigoplus_{I\in\mathcal I}\mathbb C e_I^\vee.
\]
Fix the exterior-form normalization by
\[
 (\lambda_1\wedge\cdots\wedge\lambda_n)(f_1,\ldots,f_n)
                  =\det\bigl(\lambda_i(f_j)\bigr).
\]
Consequently the continuous alternating \(n\)-forms are exactly
\(\Lambda^nD\).  Indeed, after a finite jet bound is fixed, choose an
order on its finite set of monomials.  An alternating form is the finite
sum
\[
 \sum_{I_1<\cdots<I_n}\omega(e_{I_1},\ldots,e_{I_n})
                  e_{I_1}^\vee\wedge\cdots\wedge e_{I_n}^\vee.
\]
Evaluation on ordered tuples of the finite monomial basis proves this
identity.  The determinant convention then extends it to all inputs by
multilinearity.  The ordinary scalar Chevalley--Eilenberg differential is
\begin{equation}
 \begin{split}
 (d_{\mathrm{CE}}\omega)(f_1,\ldots,f_{n+1})
  =\sum_{i<j}(-1)^{i+j}\omega\bigl(
    \{f_i,f_j\},f_1,\ldots,\widehat f_i,
                      \ldots,\widehat f_j,\ldots,f_{n+1}\bigr).
 \end{split}
 \label{eq:tb-ordinary-ce}
\end{equation}
The remaining arguments retain their order.  Scalars have trivial
\(\mathfrak h\)-action, so there are no coefficient-action terms.

If \(\omega\) uses coefficients through degree \(N\), its differential
uses coefficients through degree \(N+1\), by
\eqref{eq:tb-order-estimate}.  For every \(N\ge1\), define
\(\lambda_N(f)=[x^N]f\).  Then
\begin{equation}
 (d_{\mathrm{CE}}\lambda_N)(x^{N+1},y)=-(N+1).
 \label{eq:tb-sharp-loss}
\end{equation}
Equation~\eqref{eq:tb-sharp-loss} proves sharpness of the increase by one.
At \(N=0\), the
functional \([x^0]\) is zero on \(\mathfrak h\)'s zero-constant
representatives and \(\{x,y\}=0\).  Both sides of the asserted
nonzero phenomenon disappear.  No sharp loss is claimed for that case.
In particular, the finite coefficient projections have not been used as
Lie quotient maps.

\subsection{Suspension and signs}
\label{sec:tb-signs}

The source \(\mathcal L\) has nonzero terms in several degrees.
Suspension records the interaction of form degree and cochain degree.
Let \(L\) be a cohomologically graded differential graded Lie algebra.
Use
\[
 (L[1])^j=L^{j+1},\qquad s:L\longrightarrow L[1],\qquad |s|=-1.
\]
Put \(V=L[1]\).  The graded symmetric power \(S^n(V)\) is the
quotient of \(V^{\otimes n}\) imposing
\(vw=(-1)^{|v||w|}wv\) for homogeneous inputs.  In particular, the
square of an odd homogeneous vector is zero over \(\mathbb C\).
The empty word is \(1\in S^0(V)=\mathbb C\), and
\(S^c(V)=\bigoplus_{n\ge0}S^n(V)\).  Every element of this direct
sum has finitely many nonzero arities.  Concatenation of words gives
the graded-commutative multiplication \(\mu\) on \(S^c(V)\).

For a permutation of homogeneous inputs, its Koszul sign is the product
of \((-1)^{|v_i||v_j|}\) over pairs whose relative order is reversed.
Write \(v_S\) for the subword indexed by \(S\subset\{1,\ldots,n\}\)
in increasing order, with \(v_\varnothing=1\).  Let
\(\epsilon(S,T)\) be the Koszul sign for listing \(S\) first and
then its complement \(T\).  The unshuffle coproduct is
\begin{equation}
 \Delta(v_1\cdots v_n)=
   \sum_{S\sqcup T=\{1,\ldots,n\}}
                      \epsilon(S,T)v_S\otimes v_T.
 \label{eq:tb-coalgebra-coproduct}
\end{equation}
The coaugmentation is the inclusion \(\mathbb C\to S^0(V)\subset
S^c(V)\), and the counit is projection onto \(S^0(V)\).
For homogeneous linear maps, the tensor convention is
\begin{equation}
 (A\otimes B)(u\otimes v)=(-1)^{|B||u|}Au\otimes Bv.
 \label{eq:tb-tensor-sign}
\end{equation}

Both iterated coproducts list all ordered partitions of the inputs into
three subwords.  In either iteration, the sign is the product for the
inversions that place those subwords in the same final order.
This proves coassociativity.  Interchanging the two blocks in
\eqref{eq:tb-coalgebra-coproduct}, with the sign of their total degrees,
proves graded cocommutativity.  The empty-word terms verify the counit
and coaugmentation identities.

For an even vector \(v\), the normalization gives
\(\Delta(v^n)=\sum_{a=0}^n\binom na v^a\otimes v^{n-a}\).
For homogeneous scalar forms on the coalgebra, define their product by
\(\alpha\beta=(\alpha\otimes\beta)\Delta\), with the tensor sign
in~\eqref{eq:tb-tensor-sign}.  A homogeneous coordinate \(u\) that
vanishes on arities other than one and satisfies \(u(v)=1\) therefore
has \(u^n(v^n)=n!\).  Induction multiplies by \(n\) at the last step
of the coproduct.  This fixes the multiplicities of repeated even inputs;
the formulas do not use divided powers.

A homogeneous linear map \(A:S^c(V)\to S^c(V)\) is a coderivation
if it satisfies
\[
 \Delta A=(A\otimes\mathrm{id}+\mathrm{id}\otimes A)\Delta,
\]
where the second summand uses~\eqref{eq:tb-tensor-sign}.

\begin{lemma}[Primitives and uniqueness of coderivations]
\label{lem:tb-coderivation-uniqueness}
The elements \(z\in S^c(V)\) satisfying
\(\Delta z=1\otimes z+z\otimes1\) are precisely \(V=S^1(V)\).
Such elements are called primitive.  A coderivation that vanishes on
\(1\) is uniquely determined by its composition with the projection
\(\pi:S^c(V)\to V\), called its corestriction.
\end{lemma}
\begin{proof}
The coproduct preserves total arity, so the primitive condition can be
checked separately in each arity.
For arity \(m\ge2\), consider the component
\(\Delta_{1,m-1}:S^m(V)\to V\otimes S^{m-1}(V)\).
For each selected input, the coproduct first moves it to the front.
Multiplication moves it back to the original position with the same
Koszul sign.  Their product is one.  The \(m\) choices therefore give
\[
 \mu\Delta_{1,m-1}=m\,\mathrm{id}_{S^m(V)}.
\]
Since \(m\) is invertible in \(\mathbb C\), a primitive has no
component in arity \(m\ge2\).  A scalar \(a1\) is primitive only
if \(a1\otimes1=2a1\otimes1\), so \(a=0\).
All vectors in \(V\) are primitive by the coproduct formula.

Let \(A\) be the difference of two coderivations with the same
corestriction and with \(A1=0\).  Induct on input arity.  Suppose
\(A\) vanishes on all shorter words.  The coderivation identity,
applied to a word \(w\) of the next arity, leaves only
\[
 \Delta(Aw)=Aw\otimes1+1\otimes Aw.
\]
All other terms contain \(A\) applied to a shorter word and vanish.
Thus \(Aw\) is primitive.  Its projection to \(V\) is zero by the
corestriction hypothesis.  The first part of the proof gives \(Aw=0\).
This proves the induction and uniqueness.
\end{proof}

Define a degree-one coderivation \(Q\) by its components
\begin{equation}
 q_1(sx)=-s(dx),\qquad
 q_2(sx\,sy)=(-1)^{|x|}s[x,y],\qquad q_n=0\ (n>2).
 \label{eq:tb-q-conventions}
\end{equation}
Its arity-zero component is zero, so \(Q\) vanishes on the scalar
coaugmentation.
Here a coderivation means that applying the coproduct after \(Q\)
equals applying \(Q\) to either tensor factor after the coproduct,
with the degree-one tensor sign.  Its extension is explicit.  For
\(v_i=sx_i\), sum
\[
 \sum_i\epsilon(i;\widehat i)\,q_1(v_i)v_{\widehat i}
 \; +\sum_{i<j}\epsilon(i,j;\widehat{i,j})\,
                q_2(v_i,v_j)v_{\widehat{i,j}},
\]
where each remaining list retains its order and \(\epsilon\) is the
Koszul sign for moving the selected inputs to the front.
Each fixed arity involves a finite sum.
Write \(Q=Q_1+Q_2\) for the unary and binary extensions in this formula.

These formulas do define a coderivation.  In a coproduct of a summand,
the selected block of one or two inputs, after application of \(q_i\),
lies entirely in one output factor.  Such choices correspond exactly
to applying \(Q\) in that factor after splitting the original word.
Both orders move the same original inputs into the same ordered blocks.
Moving a degree-one map past the left output block contributes the
additional sign in~\eqref{eq:tb-tensor-sign}.  Hence their terms and
signs agree.  The formula for \(q_2\) respects graded symmetry because
the bracket is graded antisymmetric.  The uniqueness in
Lemma~\ref{lem:tb-coderivation-uniqueness} identifies this extension
with the coderivation specified by~\eqref{eq:tb-q-conventions}.

\begin{lemma}[The square-zero identity]
\label{lem:tb-q-square}
The coderivation~\eqref{eq:tb-q-conventions} satisfies \(Q^2=0\).
\end{lemma}
\begin{proof}
The square of an odd coderivation is a coderivation: in the double
coproduct calculation, the two terms applying \(Q\) in different
tensor factors cancel.  Also \(Q^2(1)=0\).
Lemma~\ref{lem:tb-coderivation-uniqueness} shows that it suffices to
compute the projection to \(L[1]\).  Since each application lowers
arity by at most one, the
only possible input arities are one, two, and three.

In arity one the expression is \(sd^2x=0\).
For \(a=|x|\), the arity-two expression is
\[
 (-1)^{a+1}s\bigl(d[x,y]-[dx,y]-(-1)^a[x,dy]\bigr)=0.
\]
For \(a=|x|\), \(b=|y|\), and \(c=|z|\), the arity-three
expression is
\[
 (-1)^b s\bigl([[x,y],z]-(-1)^{bc}[[x,z],y]
                +(-1)^{a(b+c)}[[y,z],x]\bigr).
\]
Antisymmetry changes the last term inside parentheses to
\(-[x,[y,z]]\).  The remaining equality is the graded Jacobi identity.
Thus every corestriction component vanishes, and so does \(Q^2\).
\end{proof}

For \(L=\mathfrak h\) with its usual coefficient topology, or
\(L=\mathcal L\) with its smooth topology, let
\[
 M_n^p(L)=\operatorname{Hom}_{\mathrm{cont}}^p
                        (S^n(L[1]),\mathbb C).
\]
This notation means jointly continuous graded-symmetric multilinear
forms, defined on the underlying locally convex spaces with usual scalar
values.  A homogeneous form of degree \(p\) vanishes unless its input
degrees sum to \(-p\).  No tensor completion is implicit in this notation.
Define the arity-completed cochain spaces degree by degree:
\begin{equation}
 \widehat C^p_{\mathrm{CE,cont}}(L)=\prod_{n\ge0}M_n^p(L).
 \label{eq:tb-arity-completion}
\end{equation}
Here \(M_0^0=\mathbb C\) and \(M_0^{p\ne0}=0\).

For a homogeneous cochain \(\phi\) of degree \(p\), set
\(\phi_{-1}=0\) and define
\begin{equation}
 D\phi=-(-1)^p\phi Q,\qquad
 (D\phi)_n=-(-1)^p(\phi_nQ_1+\phi_{n-1}Q_2).
 \label{eq:tb-ce-differential}
\end{equation}
The first part preserves arity, and the second raises it by one.
Both raise cohomological degree by one.  Lemma~\ref{lem:tb-q-square}
implies \(D^2=0\).

Multiplication is the graded dual of the coproduct.  For forms of specified
arities, its formula on homogeneous inputs is
\begin{equation}
 (\phi\psi)(v_1\cdots v_n)=
 \sum_{S\sqcup T=\{1,\ldots,n\}}
 \epsilon(S,T)(-1)^{|\psi|\sum_{i\in S}|v_i|}
                       \phi(v_S)\psi(v_T).
 \label{eq:tb-shuffle-product}
\end{equation}
Only sublists of the prescribed sizes contribute.  For arity products,
sum also over the finitely many pairs of arities adding to \(n\).
The tensor sign in~\eqref{eq:tb-shuffle-product} is part of the dual
pairing.  The coproduct is coassociative and graded cocommutative, so this
product is associative and graded commutative.  Applying the coderivation
identity for \(Q\) and the sign in~\eqref{eq:tb-ce-differential} gives
\(D(\phi\psi)=(D\phi)\psi+(-1)^{|\phi|}\phi(D\psi)\).

\begin{lemma}[Ordinary alternating forms]
\label{lem:tb-unsuspension}
For \(L=\mathfrak h\) in degree zero, arity \(n\) has cohomological
degree \(n\).  The identification with ordinary alternating forms is
\begin{equation}
 \omega(x_1,\ldots,x_n)
   =(-1)^{n(n-1)/2}\phi(sx_1\cdots sx_n).
 \label{eq:tb-unsuspension}
\end{equation}
It identifies the product with the ordinary exterior product and
\(D\) with~\eqref{eq:tb-ordinary-ce}.
\end{lemma}
\begin{proof}
Every suspended input has degree \(-1\), so graded symmetry is ordinary
antisymmetry.  Put \(t_n=(-1)^{n(n-1)/2}\).
Moving two selected inputs \(i<j\) to the front gives sign
\((-1)^{i+j-3}\).  The ratio of unsuspension signs is
\(t_{n+1}/t_n=(-1)^n\).  Multiplication by the differential sign
\(-(-1)^n\) yields the coefficient \((-1)^{i+j}\) of
\eqref{eq:tb-ordinary-ce}.

For input arities \(a,b\), the extra tensor evaluation sign in
\eqref{eq:tb-shuffle-product} is \((-1)^{ab}\).
The equality \(t_{a+b}=(-1)^{ab}t_at_b\) cancels this extra sign
under~\eqref{eq:tb-unsuspension}, leaving the exterior shuffle product.
\end{proof}

For coordinates \(c^I(se_J)=\delta^I_J\), write
\(\{e_I,e_J\}=\sum_K C^K_{IJ}e_K\).  The preceding signs give
\[
 Dc^K=-\tfrac12\sum_{I,J}C^K_{IJ}c^Ic^J.
\]
For each fixed \(K\), the sum is finite: nonzero terms have
\(|I|+|J|=|K|+2\), with both input degrees positive.
The factor \(1/2\) counts the two orders of each pair.

\subsection{Bounded products and degreewise completeness}
\label{sec:tb-completeness}

A subset \(B\) of a locally convex space is bounded if every continuous
seminorm is bounded on \(B\).  On continuous \(n\)-linear forms use
the topology with seminorms
\begin{equation}
 p_{B_1,\ldots,B_n}(\omega)
    =\sup_{x_i\in B_i}|\omega(x_1,\ldots,x_n)|,
 \label{eq:tb-bounded-topology}
\end{equation}
where all \(B_i\) are bounded.  These seminorms are finite.  A continuous
multilinear form is bounded on a product of sufficiently small
zero-neighborhoods, and bounded sets can be scaled into those
neighborhoods.  Give~\eqref{eq:tb-arity-completion} the product topology
of~\eqref{eq:tb-bounded-topology} in each fixed degree.

\begin{lemma}[Completeness of continuous multilinear forms]
\label{lem:tb-bounded-complete}
If \(E\) is a metrizable locally convex complex space, the continuous
scalar \(n\)-linear forms on \(E^n\) are complete for
\eqref{eq:tb-bounded-topology}.  Closed symmetry and homogeneous-degree
conditions preserve completeness.
\end{lemma}
\begin{proof}
Arity zero is the complete space \(\mathbb C\).  Suppose \(n\ge1\),
and let \((\omega_a)\) be a Cauchy net of continuous forms.
Singletons are bounded, so the net has a pointwise limit \(\omega\).
The defining multilinear identities pass to pointwise limits.

Fix bounded sets \(B_i\) and \(\epsilon>0\).  There is an index
\(a_0\) such that
\[
 p_{B_1,\ldots,B_n}(\omega_a-\omega_b)<\epsilon
                         \qquad(a,b\ge a_0).
\]
Letting \(b\) tend pointwise to the limit gives
\(p_{B_1,\ldots,B_n}(\omega_a-\omega)\le\epsilon\) for
\(a\ge a_0\).  One continuous member \(\omega_a\) is bounded on
this product, so \(\omega\) is bounded there as well.  Thus convergence
is uniform on every product of bounded sets.

Continuity follows if a multilinear form is bounded on a product
\(V_1\times\cdots\times V_n\) of open balanced convex neighborhoods
of zero.  Their continuous Minkowski seminorms are
\[
 p_i(x)=\inf\{t>0:x\in tV_i\}.
\]
Here balanced means that \(aV_i\subset V_i\) for \(|a|\le1\),
and convex means that \(tx+(1-t)y\in V_i\) for \(x,y\in V_i\)
and \(0\le t\le1\).  Finite intersections of seminorm balls give
such neighborhoods.  They absorb each vector because sufficiently small
scalar multiples approach zero, so \(p_i\) is finite.
Balancedness gives absolute homogeneity of \(p_i\).  Convexity gives
subadditivity: if \(x\in aV_i\) and \(y\in bV_i\), then
\(x+y\in(a+b)V_i\), and taking infima gives the inequality.

Openness gives \(\{p_i<\epsilon\}=\epsilon V_i\).
For one inclusion, \(p_i(x)<\epsilon\) puts \(x\) in some
\(tV_i\) with \(t<\epsilon\).  Balancedness then gives
\(x\in\epsilon V_i\).  Conversely, if \(x/\epsilon\in V_i\),
openness along the real ray gives \((1+\delta)x/\epsilon\in V_i\)
for some \(\delta>0\), so \(p_i(x)<\epsilon\).
These open balls prove continuity at zero.  The inequality
\(|p_i(x)-p_i(y)|\le p_i(x-y)\) proves continuity everywhere.
Scaling each input gives a bound
\[
 |\omega(x_1,\ldots,x_n)|\le C\prod_i p_i(x_i).
\]
The zero-seminorm cases follow by a limit.  This bound proves continuity
at the origin and, after expanding increments, at arbitrary input tuples.

If \(\omega\) were discontinuous, it would therefore be unbounded
on every product of zero-neighborhoods.  Choose a decreasing countable
neighborhood base \(U_k\) at zero.  There would be vectors
\(x_{i,k}\in U_k\) with
\[
 |\omega(x_{1,k},\ldots,x_{n,k})|>k.
\]
Each sequence \(x_{i,k}\to0\) is bounded: every continuous seminorm
has a bounded tail and finitely many remaining values.  Boundedness of
\(\omega\) on the product of these sequences contradicts the displayed
inequality.  Thus \(\omega\) is continuous, proving completeness.
Symmetry and degree restrictions are pointwise equalities and hence
define closed subspaces.
\end{proof}

\begin{proposition}[The two degreewise completions]
\label{prop:tb-complete}
For \(L=\mathfrak h\) with usual coefficients or \(L=\mathcal L\),
every \(\widehat C^p_{\mathrm{CE,cont}}(L)\) is complete for the
bounded-product topology.  It is also complete for the arity filtration
\[
 F^r\widehat C^p=\prod_{n\ge r}M_n^p(L),\qquad
 \widehat C^p\cong\varprojlim_r\widehat C^p/F^r\widehat C^p.
\]
The differential and multiplication are continuous degreewise, and
\(D F^r\subset F^r\), \(F^rF^s\subset F^{r+s}\).
\end{proposition}
\begin{proof}
The finite graded sums underlying \(\mathfrak h[1]\) and
\(\mathcal L[1]\) are metrizable locally convex spaces.
Lemma~\ref{lem:tb-bounded-complete} proves completeness of each
\(M_n^p(L)\).  A Cauchy net in their product converges in each
coordinate, and these limits give its limit in the product.  This proves
bounded-product completeness.

The quotient by \(F^r\) retains precisely the arities below \(r\).
A compatible family of these finite lists specifies one element of
the full arity product, uniquely.  This proves the inverse-limit
statement independently of the preceding topological completeness proof.
The arity formulas for multiplication and differential give the stated
filtration relations.

For bounded sets, a continuous linear map sends a bounded set to a
bounded set.  A continuous bilinear map sends a product of bounded sets
to a bounded set: apply its seminorm estimates and boundedness of each
input set.  Thus precomposition with \(d\) and with the bracket is
continuous for~\eqref{eq:tb-bounded-topology}.  Each differential
component is a finite sum of such precompositions.
Every shuffle term is bounded by the product of the two corresponding
seminorms of its factors.  This proves joint continuity of multiplication
at specified arities.  A fixed output arity uses only finitely many
input arities, proving degreewise continuity on the product.
\end{proof}

The word complete here always applies in a fixed cohomological degree.
For \(\mathfrak h\), only arity \(p\) occurs in degree \(p\).
For \(\mathcal L\), suspended degrees \(-1,0,1\) allow arbitrarily
large arities in a fixed degree, so the arity product is substantive.
No product over all cohomological degrees is taken.

An ungraded direct sum of degrees with the arity-product subspace topology
would not be complete.  Choose distinct coordinates \(c^1,c^2,\ldots\)
of \(D\).  The sequence
\[
 \sum_{n=1}^N c^1c^2\cdots c^n
\]
converges in each arity as \(N\to\infty\), but its limit has nonzero
components in infinitely many cohomological degrees.  That limit does
not lie in the ungraded direct sum.  Adding it would require a separate
total completion, outside~\eqref{eq:tb-arity-completion}.

\section{Taylor pullback of scalar cochains}
\label{sec:tb-pullback}

Finite holomorphic derivative expressions pull back along Taylor
extraction even though arbitrary formal cocycles have no holomorphic
lifts.  The resulting map is contravariant, as a scalar functional is
evaluated by precomposing its arguments.

\begin{proposition}[The scalar cochain map]
\label{prop:tb-eight}
For usual complex coefficient topologies, Taylor pullback gives an
injective continuous morphism of differential graded algebras with
graded-commutative products, complete in every cohomological degree,
\begin{equation}
 \bar J^*:\widehat C^\bullet_{\mathrm{CE,cont}}(\mathfrak h)
       \longrightarrow\widehat C^\bullet_{\mathrm{CE,cont}}(\mathcal L).
 \label{eq:tb-ce-pullback}
\end{equation}
It preserves arity and cohomological degree.  Its image consists of
cochains that vanish on every tuple with a positive Dolbeault-degree
argument and, in each arity, depend on finitely many positive-order pure
holomorphic derivatives at zero.
\end{proposition}
\begin{proof}
For each arity, define
\[
 (\bar J^*\phi)_n=\phi_n\circ S^n(\bar J[1]).
\]
Proposition~\ref{prop:tb-one} gives the exact identity
\[
 S^c(\bar J[1])Q_{\mathcal L}
        =Q_{\mathfrak h}S^c(\bar J[1]).
\]
Precomposition therefore commutes with
\eqref{eq:tb-ce-differential}.  The map on symmetric coalgebras preserves
the coproduct, so it also preserves~\eqref{eq:tb-shuffle-product}.
The induced map \(\bar J[1]\) has degree zero, which proves the grading
assertions.

Continuity follows first at the level of multilinear forms from
Proposition~\ref{prop:tb-six}.  For the bounded-product topology it is
given exactly by
\[
 p_{B_1,\ldots,B_n}(\bar J^*\phi)
   =p_{\bar J[1](B_1),\ldots,\bar J[1](B_n)}(\phi).
\]
The images are bounded because \(\bar J[1]\) is continuous linear.
This proves continuity in each arity and then on the arity product.

If a pulled-back cochain vanishes, choose smooth lifts of every member
of an arbitrary finite tuple of formal inputs by
Proposition~\ref{prop:tb-two}.  Evaluation on this tuple proves that
the original cochain vanishes.  No linear or continuous choice of lifts
is needed.  Thus the map is injective.

Proposition~\ref{prop:tb-seven} gives finite jet support of formal scalar
cochains.  Taylor extraction kills positive Dolbeault degrees and retains
only pure holomorphic derivatives.  The zeroth derivative disappears
because the source is already modulo constants.
Conversely, any multilinear cochain with the stated properties factors
through a finite positive-order holomorphic jet projection in each
argument.  These projections are surjective by
Proposition~\ref{prop:tb-two}, so the factor is a unique scalar form
on the corresponding formal coefficient spaces.  It is continuous by
Proposition~\ref{prop:tb-seven}.  This proves the image description.
\end{proof}

The internal differential of a pulled-back scalar cochain vanishes because
it would insert a positive-degree Dolbeault argument.  Its remaining
differential is precisely the Hamiltonian bracket differential
\eqref{eq:tb-ordinary-ce}.  This is the compatibility needed to evaluate
formal scalar symmetry cochains on smooth parameters at the origin.

For the adic formal space, Proposition~\ref{prop:tb-seven} gives the same
finite jet formulas as vector spaces.  Composing these formulas with
smooth jets still gives continuous cochains with usual complex values.
This construction is not the application of a continuous Lie map into
the adic target: Proposition~\ref{prop:tb-six} excludes that continuity.
It is an explicit algebraic identification of finite jet forms followed
by the usual-valued pullback~\eqref{eq:tb-ce-pullback}.

Discrete scalar values on the smooth side give a different answer.
A positive-arity continuous multilinear map to \(\mathbb C_{\mathrm{disc}}\)
vanishes.  Fix all but one argument.  The remaining continuous linear map
has connected domain and discrete image, so its image is the singleton
zero.  Thus the scalar topology cannot be omitted from the comparison.

Injectivity in~\eqref{eq:tb-ce-pullback} is injectivity on cochains.
It does not prove injectivity on Chevalley--Eilenberg cohomology: a
primitive in the larger complex need not belong to the image.
It also gives no coefficient-valued or cotangent comparison.
The representation and dual calculations below locate two separate
obstructions to such extensions.

\section{Translations, representations, and dual contraction}
\label{sec:tb-duality}

\begin{proposition}[Finite-dimensional representations of the full algebra]
\label{prop:tb-nine}
Every continuous complex-linear Lie homomorphism
\(\rho:\mathfrak h\to\operatorname{End}(V)\), with \(V\)
finite-dimensional, is zero.  This holds for either formal source
topology and for either usual or discrete finite-dimensional target
topology.  The algebra \(\mathfrak h\) has no proper open Lie ideal
and is perfect.
\end{proposition}
\begin{proof}
For a discrete target, continuity gives a tail
\(\mathfrak m^N\subset\ker\rho\).  For a usual finite-dimensional
target, the inverse image of a bounded norm ball contains a tail subspace.
Scaling that subspace forces its image to be zero, giving the same
inclusion.  The kernel is a Lie ideal.

For arbitrary \(f=\sum_{a+b\ge1}f_{ab}x^ay^b\), define
\begin{equation}
 I_y^Nf=\sum_{a+b\ge1}
       \frac{b!}{(b+N)!}f_{ab}x^ay^{b+N}.
 \label{eq:tb-antiderivative}
\end{equation}
It lies in \(\mathfrak m^{N+1}\subset\ker\rho\), and
\((\operatorname{ad}_x)^NI_y^Nf=f\).  Every intermediate term has
positive degree, so removal of constants causes no change.
Ideal stability gives \(f\in\ker\rho\).

Every open Lie ideal contains a tail, so the same argument makes it all
of \(\mathfrak h\).  Taking one antiderivative gives
\(f=\{x,I_yf\}\).  Thus \([\mathfrak h,\mathfrak h]=\mathfrak h\).
\end{proof}

For trivial scalar coefficients, a one-cocycle vanishes on all brackets.
Perfectness therefore gives
\(H^1_{\mathrm{CE}}(\mathfrak h,\mathbb C)=0\), even without
continuity restrictions.

\begin{example}[The origin-preserving algebra]
\label{ex:tb-origin}
The distinct subalgebra \(\mathfrak h_+=\mathfrak m^2\) preserves
the origin.  Estimate~\eqref{eq:tb-order-estimate} makes
\(\mathfrak m^{N+1}\) a Lie ideal in \(\mathfrak h_+\) for
\(N\ge1\).  Its action on linear functions is
\[
 \rho_+:\mathfrak h_+\longrightarrow
       \operatorname{End}(\mathfrak m/\mathfrak m^2),
 \qquad \rho_+(f)[\ell]=[\{f,\ell\}].
\]
The bracket derivation preserves \(\mathfrak m\) and \(\mathfrak m^2\),
so the action is defined.  Jacobi proves that it is a Lie representation.
Its kernel is \(\mathfrak m^3\).  Indeed, higher terms act trivially,
and both partial derivatives of a quadratic vanish only for the zero
quadratic.

With \(e=x^2/2\), \(f=-y^2/2\), and \(h=-xy\), the matrices in
the ordered basis \((x,y)\) are
\[
 \rho_+(e)=\begin{pmatrix}0&1\\0&0\end{pmatrix},\quad
 \rho_+(f)=\begin{pmatrix}0&0\\1&0\end{pmatrix},\quad
 \rho_+(h)=\begin{pmatrix}1&0\\0&-1\end{pmatrix}.
\]
Direct Poisson calculation gives
\(\{h,e\}=2e\), \(\{h,f\}=-2f\), and \(\{e,f\}=h\).
Thus \(\mathfrak m^2/\mathfrak m^3\cong\mathfrak{sl}_2\).
\end{example}

The nonzero representation in Example~\ref{ex:tb-origin} is continuous
from usual coefficients to the usual matrix topology, since it depends
on only three coefficients.  From the adic topology it is continuous
to either usual or discrete matrices, since its kernel contains the open
tail \(\mathfrak m^3\).  From usual coefficients to nonzero discrete
matrices it is discontinuous: the source is connected, and a continuous
map from it to a discrete space has constant image.  These topology
choices are part of the representation statement.
The translation identity \(\{x,y^2\}=2y\) proves that
\(\mathfrak h_+\) is not an ideal of \(\mathfrak h\).

\subsection{The continuous dual}
\label{sec:tb-dual}

The vector space \(D=\bigoplus_I\mathbb C e_I^\vee\) from
Proposition~\ref{prop:tb-seven} admits two relevant topologies.
For adic \(\mathfrak h\) over \(\mathbb C_{\mathrm{disc}}\), give
\(D\) the discrete topology.  For the usual coefficient product,
give \(D\) its usual locally convex direct-sum topology, denoted
\(D_{\mathrm{lc}}\).  This topology is generated by
\begin{equation}
 p_a(\lambda)=\sum_Ia_I|\lambda_I|,
                  \qquad a_I\ge0\text{ arbitrary}.
 \label{eq:tb-dual-seminorms}
\end{equation}
Each sum is finite on a given dual vector.  Every seminorm \(p\) on
the algebraic direct sum is bounded above by the one with
\(a_I=p(e_I^\vee)\).  Thus this is the finest locally convex topology
on that direct sum.

\begin{proposition}[Evaluation in the two categories]
\label{prop:tb-ten}
The evaluation map
\[
 \mathfrak h_{\mathrm{adic}}\times D_{\mathrm{disc}}
       \longrightarrow\mathbb C_{\mathrm{disc}},\qquad
 (f,\lambda)\longmapsto\lambda(f)
\]
is jointly continuous.  The evaluation map
\[
 \mathfrak h_{\mathrm{prod}}\times D_{\mathrm{lc}}
       \longrightarrow\mathbb C
\]
is separately continuous and is not jointly continuous.
\end{proposition}
\begin{proof}
In the adic case, fix \((f_0,\lambda_0)\), and choose \(N\) containing
the support of \(\lambda_0\).  Evaluation is constant on the open set
\[
 (f_0+\mathfrak m^{N+1})\times\{\lambda_0\}.
\]
This proves joint continuity at every point.

In the usual category, a fixed \(\lambda\) uses finitely many
coefficient evaluations, so evaluation is continuous in \(f\).
For fixed \(f=(f_I)\), the estimate
\(|\lambda(f)|\le p_{|f|}(\lambda)\) proves continuity in \(\lambda\).

Suppose joint continuity held at zero.  There would be neighborhoods
\(U_0\) and \(V_0\) whose pairing lies in \(|z|<1/2\).
The neighborhood \(U_0\) contains a basic product neighborhood
restricting finitely many coordinates.  Choose an unrestricted index
\(I\).  Then \(te_I\in U_0\) for every \(t\in\mathbb C\).
The coordinate-line inclusion \(\mathbb C e_I^\vee\to D_{\mathrm{lc}}\)
is continuous.  Hence \(V_0\) contains \(se_I^\vee\) for some
\(s\ne0\).  Taking \(t=1/s\) gives pairing one, a contradiction.
\end{proof}

The topology \(D_{\mathrm{lc}}\) is also the strong continuous-dual
topology of \(\mathfrak h_{\mathrm{prod}}\).  A subset of the product
is bounded exactly when every coordinate is bounded.  Indeed, the unit
neighborhood of any continuous seminorm contains a full coefficient tail.
Scaling shows that the seminorm vanishes on that tail.  It therefore
factors through a finite coordinate space, where coordinatewise bounded
sets are bounded.  The converse follows by taking individual coordinate
seminorms.  For arbitrary finite-valued
weights \(a_I\ge0\), the rectangle
\[
 B_a=\{f:|f_I|\le a_I\text{ for all }I\}
\]
is bounded, and
\[
 \sup_{f\in B_a}|\lambda(f)|=\sum_Ia_I|\lambda_I|.
\]
Conversely, any bounded set is contained in such a rectangle.
Thus uniform convergence on bounded subsets gives exactly
\eqref{eq:tb-dual-seminorms}.

\begin{proposition}[The coadjoint continuity obstruction]
\label{prop:tb-coadjoint}
The coadjoint action
\[
 (f\cdot\lambda)(g)=-\lambda(\{f,g\})
\]
is jointly continuous from
\(\mathfrak h_{\mathrm{adic}}\times D_{\mathrm{disc}}\) to
\(D_{\mathrm{disc}}\).  It is separately continuous but not jointly
continuous from \(\mathfrak h_{\mathrm{prod}}\times D_{\mathrm{lc}}\)
to \(D_{\mathrm{lc}}\).
\end{proposition}
\begin{proof}
If \(\lambda\) is supported through degree \(N\), its coadjoint
image is supported through degree \(N+1\) and depends only on \(f\)
through degree \(N+1\).  This follows from
\eqref{eq:tb-order-estimate}.  At a fixed \(\lambda\), the map in
\(f\) is therefore continuous in both stated categories.  In the
adic category, holding \(\lambda\) fixed in an open singleton proves
joint continuity.

For fixed \(f\) in the usual category, the action is a linear map
from a space with the finest locally convex topology.  The pullback
of each target seminorm is a seminorm and hence continuous, so this
linear map is continuous.  This proves separate continuity.

The formal antiderivative \(I_x\), defined as in
\eqref{eq:tb-antiderivative} with the variables exchanged, is continuous
for usual coefficients.  If the usual coadjoint action were jointly
continuous, composition with evaluation at the fixed vector \(y\)
would make the following map jointly continuous:
\[
 (f,\lambda)\longmapsto(I_xf\cdot\lambda)(y)
       =-\lambda(\{I_xf,y\})=-\lambda(f).
\]
This contradicts Proposition~\ref{prop:tb-ten}.
\end{proof}

The coadjoint formula is a Lie action because Jacobi gives
\[
 [f,g]\cdot\lambda=f\cdot(g\cdot\lambda)-g\cdot(f\cdot\lambda).
\]
Evaluation on any third Hamiltonian reduces this equation to its Jacobi
identity.  The algebraic coadjoint semidirect product therefore has a continuous
Lie bracket in the adic/discrete category.  With the usual product and
locally convex direct-sum topologies, that bracket fails joint continuity.
The two categories cannot be identified through their equal underlying
dual vector spaces.

\subsection{The cotangent boundary of the scalar comparison}
\label{sec:tb-cotangent-boundary}

In the discrete formal category, consider the graded Lie algebra
\(\mathfrak h\ltimes D[1]\), where \(D[1]\) is abelian and
\(\mathfrak h\) acts by the coadjoint action.
Its underlying graded space has \(\mathfrak h\) in degree zero and
\(D[1]\) in degree \(-1\), with zero differential.  Writing
\(s\lambda\) for a vector of \(D[1]\), the brackets are
\[
 [f,g]=\{f,g\},\qquad
 [f,s\lambda]=s(f\cdot\lambda),\qquad
 [s\lambda,s\mu]=0.
\]
The reverse cross bracket is fixed by graded antisymmetry.
The Jacobi identity with two Hamiltonians and one dual input is the
coadjoint action identity proved above.  Identities with at least two
dual inputs vanish because the dual summand is abelian.  This verifies
the graded Lie structure on the stated carrier.
Under the suspension convention~\eqref{eq:tb-q-conventions}, the
Chevalley--Eilenberg coordinates dual to \(s\mathfrak h\) and
\(s^2D\) have degrees
\[
 |c^I|=1,\qquad |u_I|=2.
\]
The coordinate contraction specified by
\([c^I,u_J]=\delta^I_J\) has degree \(-3\).
Giving \(u_I\) weight one and \(c^I\) weight zero, then subtracting
twice this weight from the degree, gives \(|u_I|=0\), \(|c^I|=1\),
and contraction degree \(-1\).  These are grading calculations on
the indicated linear coordinates.

Suppose a graded locally convex space \(\mathcal B\) has a
distinguished scalar copy \(\mathbb C\cdot1\subset\mathcal B^0\)
and continuous linear inclusions
\[
 i_D:D_{\mathrm{lc}}\longrightarrow\mathcal B^1,
 \qquad i_h:\mathfrak h_{\mathrm{prod}}\longrightarrow\mathcal B^2,
 \qquad i_D(e_I^\vee)=c^I,\quad i_h(e_J)=u_J,
\]
and a continuous linear projection \(p:\mathcal B^0\to\mathbb C\)
with \(p(1)=1\).  Assume a jointly continuous bilinear contraction
of degree \(-3\) has the displayed coordinate values.
The resulting scalar pairing is
\[
 B(\lambda,F)=p\bigl([i_D(\lambda),i_h(F)]\bigr).
\]
For a finite polynomial \(P\) without constant term, the coordinate
identities and finite multilinearity give \(B(\lambda,P)=\lambda(P)\),
since \(\lambda\) also has finite support.
For arbitrary \(F\in\mathfrak h\), let \(P_N\) be its polynomial
truncation through degree \(N\).  Then \(P_N\to F\) in the usual
coefficient topology.  Continuity at each fixed \(\lambda\) gives
\[
 B(\lambda,F)=\lim_N B(\lambda,P_N)
            =\lim_N\lambda(P_N)=\lambda(F).
\]
Thus joint continuity of the contraction would make the full evaluation
pairing jointly continuous, contradicting Proposition~\ref{prop:tb-ten}.
This argument uses density of finite polynomials rather than an
operation on an unspecified infinite sum of coordinates.
After the stated regrading, the same proof has \(i_h\) in degree zero
and the contraction in degree \(-1\).
Proposition~\ref{prop:tb-coadjoint} already prevents the corresponding
usual coadjoint semidirect product from being a jointly continuous
topological Lie algebra.

The scalar map~\eqref{eq:tb-ce-pullback} is unaffected by this obstruction.
It evaluates finite holomorphic jets, preserves the full Hamiltonian
bracket including translations, and respects the Dolbeault differential.
The underlying Dolbeault-to-formal map \(\bar J\) still has the kernel
cohomology \(R/R_U\) in~\eqref{eq:tb-kernel-cohomology}.
An extension to coefficient-valued or cotangent cochains requires new
carrier and continuity data.  The dual of \(\bar J:\mathcal L\to
\mathfrak h\) goes from \(D\) into the smooth continuous dual, so a
same-direction cotangent map also requires a separate construction.
Neither an analytic equivalence nor a cotangent comparison follows from
the proved scalar cochain map.
